% concatenated from The_groups_ozone_the_algebra_U_q_+_B_2_.tex
\documentclass[11pt]{article}
\usepackage{amsmath,amscd}
\usepackage{amssymb,latexsym,amsthm}
\usepackage{color,palatino}
%\usepackage[spanish,activeacute]{babel}
\usepackage[utf8]{inputenc}
\usepackage[T1]{fontenc}
\usepackage{csquotes}
\usepackage{latexsym}
\usepackage{multirow}
\usepackage{lscape}
\usepackage{enumerate}
\usepackage{longtable}
\usepackage{array}
\usepackage{amsmath}
\usepackage{amssymb}
\usepackage{fancyhdr}
\usepackage{pb-diagram}
\usepackage{amsfonts}
\usepackage{hyperref}
\usepackage{url}
\hypersetup{
  colorlinks=true,
  linkcolor=blue,
  citecolor=black,
  urlcolor=cyan
}
\usepackage{pb-diagram}
\newtheorem{theorem}{Theorem}[section]
\newtheorem{proposition}[theorem]{Proposition}
\newtheorem{lemma}[theorem]{Lemma}

%\newtheorem{conjetura}[theorem]{Conjetura}
\newtheorem{corollary}[theorem]{Corollary}

\newtheorem{examples}[theorem]{examples}
\newtheorem{definition}[theorem]{Definition}
%\newtheorem{ejercicio}[theorem]{Ejercicio}
%\newtheorem{notacion}[theorem]{Notación}
%\newtheorem{problema}[theorem]{Problema}
%\newtheorem{solucion}[theorem]{Solución}
%\newtheorem{resumen}[theorem]{Resumen}
%\newtheorem{clasificacion}[theorem]{Clasificación}
\newtheorem{question}[theorem]{Question}
\newtheorem*{question*}{Question}
\theoremstyle{definition}
\newtheorem{remark}[theorem]{Remark}
\newtheorem{note}[theorem]{Note}
\newtheorem{example}[theorem]{Example}
% Número (p,q): [n]_{p,q}
\newcommand{\pqnum}[1]{\left[#1\right]_{p,q}}
\newcommand{\Hpq}{H_{p,q}}
\newcommand{\K}{\Bbbk}
\newcommand{\ordh}{\ell}
\newcommand{\Z}{\mathbb{Z}}
\newcommand{\N}{\mathbb{N}}
\newcommand{\C}{\mathbb{C}}
\newcommand{\nuevo}[1]{#1}
\DeclareMathOperator{\supp}{supp}
\DeclareMathOperator{\GKdim}{GKdim}
\DeclareMathOperator{\Kdim}{Kdim}
\hoffset-0.2in
\voffset-2cm \setlength{\oddsidemargin}{9mm}
\setlength{\textheight}{23cm}
\setlength{\textwidth}{16cm}
\makeatletter
\def\@roman#1{\romannumeral #1}
\makeatother

\title{\textbf{The ozone group of $U_q^+(B_2)$}}
\author{James Gómez and Helbert Venegas}



\date{}
\begin{document}
\makeatletter
\def\@roman#1{\romannumeral #1}
\makeatother
\maketitle
\begin{abstract}
\noindent
We determine the ozone group of \(U_q^{+}(B_2)\). More precisely, if
\(\ell=\operatorname{ord}(q^2)\) when \(q\) is a root of unity, we obtain
$$
\operatorname{Oz}(U_q^{+}(B_2))
\cong
\begin{cases}
\mu_2 , & \text{if \(q\) is not a root of unity},\\[2mm]
\mu_{\gcd(\ell,2)}, &
\text{if \(q\) is a primitive \(m\)-th root of unity, \(m\geq5\).}
\end{cases}
$$
We also study several homological properties of \(U_q^{+}(B_2)\), showing
that it is Artin--Schelter regular of global dimension \(4\), Auslander-regular, 
Cohen--Macaulay, strongly Noetherian, and skew Calabi--Yau. Finally, in
the root-of-unity case, we relate the ozone group to the normal elements
of \(U_q^{+}(B_2)\) and show that, when \(\ell\) is odd,
\(U_q^{+}(B_2)\) is Calabi--Yau and its center is Gorenstein.
 
\bigskip

\noindent
\textit{Keywords: Quantum algebra; PI algebras; center; ozone groups.} 
\bigskip

\noindent 2020 \textit{Mathematics Subject Classification.} Primary: 16W20.
Secondary: 17B37, 16U70.
\end{abstract}
%%%%%%%%%%%%%%%%%%%%%%%%%%%%%%%%%%%%%%%%%%%%%%%%%
\section{Introduction}
The algebra \(U_q^{+}(B_2)\)  was originally defined within the general framework of quantum enveloping algebras developed by Drinfeld \cite{Drinfeld1986} and Jimbo \cite{Jimbo1985} in the 1980s, as part of the theory of quantum groups arising from the study of quantum integrable systems and the Yang-Baxter equation.
Its significance lies in the fact that it constitutes a fundamental example of a rank-$2$ quantum algebra (with two simple roots), associated to the complex simple Lie algebra $\mathfrak{g}$  of type $B_2$. In particular $U_q^{+}(B_2)$ is the positive part of the quantum group $U_q(\mathfrak{g})$, generated by the Chevalley generators $E_1$, $E_2$ corresponding to the simple roots of the root system of type $B_2$. For the case where \(q\) is not a root of unity, this algebra has been studied in detail by Launois \cite{Launois2007}, who explicitly described its prime and primitive spectra, computed its automorphism group, and obtained quantum analogues of the Weyl algebra as simple quotients. Additionally, in \cite{AndruskiewitschDumas2008} some results related to Hopf algebra automorphism are also developed. 
Recently, some properties have been studied in the case where the parameter $q$ is a root of unity. In particular, Bera and Mukherjee \cite{BeraMukherjee2026UqB2} proved that $U_q^{+}(B_2)$ is a PI algebra when $q$ is a primitive $m$-th root of unity with $m\geq 5$, computed its center, and classified its simple modules up to isomorphism. \\

Throughout this paper, \(\K\) denotes an algebraically closed field of
characteristic zero, and all algebras are assumed to be
\(\K\)-algebras.\\

\noindent Formally, the algebra   \(U_q^{+}(B_2)\) is the \(\K\)-algebra generated by two indeterminates  $e_1$ and $e_2$ subject to the quantum Serre relations,
\begin{align*}
e_1^{2}e_2 - (q^{2}+q^{-2})\,e_1 e_2 e_1 + e_2 e_1^{2} &= 0\\
e_2^{3}e_1 - (q^{2}+1+q^{-2})\,e_2^{2} e_1 e_2 + (q^{2}+1+q^{-2})\,e_2 e_1 e_2^{2}   - e_1 e_2^{3} &= 0
\end{align*}

Andruskiewitsch and Dumas \cite[Section 3.1.2]{AndruskiewitschDumas2008},   introduced a new set of generators for
$U_q^{+}(B_2)$. For this, they defined
$e_{3}=e_{1}e_{2}-q^{2}e_{2}e_{1}
$ and $
z=e_{2}e_{3}-q^{2}e_{3}e_{2}$. Thus, the set 
\(\{\, z^{i} e_{3}^{j} e_{1}^{k} e_{2}^{\ell} \mid i,j,k,\ell \in \mathbb{N}\,\}\) 
is a PBW basis for $U_q^{+}(B_2)$  and the latter turns out to be  the \(\K\)-algebra generated by
\(e_1,e_2,e_3\) and \(z\) subject to the relations:
\begin{align*}
&e_i z = z e_i, \ \text{ for }\, i=1,2,3,
&e_1 e_3 &= q^{-2} e_3 e_1,\\
&e_2 e_1 = q^{-2} e_1 e_2 - q^{-2} e_3, &e_2 e_3 &= q^{2} e_3 e_2 + z.
\end{align*}
Furthermore, \(U_q^{+}(B_2)\) can be presented as an iterated Ore extension, 
\[
U_q^{+}(B_2) \;=\; \K[z,e_3][\,e_1;\sigma_1\,][\,e_2;\sigma_2,\delta_2\,],
\]
where
\begin{align*}
\sigma_1(z) &= z, 
& \sigma_1(e_3) &= q^{-2}e_3, \\[6pt]
\sigma_2(z) &= z, 
& \sigma_2(e_3) &= q^{2}e_3, 
& \sigma_2(e_1) &= q^{-2}e_1, \\[6pt]
\delta_2(z) &= 0, 
& \delta_2(e_3) &= z, 
& \delta_2(e_1) &= -q^{-2}e_3.
\end{align*}
In 2024, Kenneth Chan, Jason Gaddis, Robert Won, and James J. Zhang introduced the notion of the \textbf{ozone group} of an algebra \(A\) (see \cite{ChanGaddisWonZhang2024OzoneGroupsCenters}), denoted by \(\operatorname{Oz}(A)\), and defined by
$$
\operatorname{Oz}(A):=\operatorname{Aut}_{Z\text{-alg}}(A).
$$
\noindent Thus, \(\operatorname{Oz}(A)\) consists of those automorphisms of \(A\) that fix its center pointwise.  If \(\sigma\in\operatorname{Oz}(A)\), then \(\sigma\) acts as the identity on the center \(Z(A)\), that is,
\(\left.\sigma\right|_{Z(A)}=\operatorname{id}_{Z(A)}\). \\

The study of ozone groups of PI Artin--Schelter regular algebras has seen significant recent development. Chan, Gaddis, Won, and Zhang
\cite{ChanGaddisWonZhang2025} established general results concerning their
finiteness, their relationship with normal elements, the rank over the center,
and connections with the Calabi--Yau and Gorenstein properties. Subsequently,
Liu, Wu, and Zhu \cite{LiuWuZhu2026} proved that the ozone group of every
noetherian PI Artin--Schelter regular algebra is abelian. In this direction,
Gaddis and Yee \cite{GaddisYee2026} recently studied the family
\(B_q(f)\), obtaining families of Calabi--Yau algebras with trivial ozone
group, while the authors \cite{GomezVenegas2026} considered the
ozone groups of this family and showed that, when \(q\) has order \(n>1\),
\(\operatorname{Oz}(B_q(f))\cong\mu_e\times\mu_e\), where
\(e=\gcd(n,\{j+1:j\in\operatorname{supp}(f)\})\). \\

For \(U_q^{+}(B_2)\), although its PI property, PI degree, and center at
roots of unity have recently been studied
\cite{BeraMukherjee2026UqB2}, an explicit computation of its ozone group
does not appear to have been previously recorded. In this work, we prove
that \(U_q^{+}(B_2)\) is Artin--Schelter regular of global dimension \(4\)
and, writing \(\ell=\operatorname{ord}(q^2)\), we obtain
\(\operatorname{Oz}(U_q^{+}(B_2))\cong\mu_{\gcd(\ell,2)}\).
Among other consequences, this yields the Calabi--Yau property of
\(U_q^{+}(B_2)\) when \(\ell\) is odd. The homological analysis is carried
out by realizing \(U_q^{+}(B_2)\) as a graded skew PBW extension, which
allows us to apply the results on Auslander regularity, the
Cohen--Macaulay property, and strong Noetherianity from
\cite{LezamaVenegas2017Homological}. We also show that, when \(\ell\) is odd, \(U_q^+(B_2)\) is Calabi--Yau and its center \(Z(U_q^+(B_2))\) is Gorenstein.


\section{Preliminaries}
In this section, we introduce the main notions and homological properties that will be used throughout the paper.\\

We recall that a \(\K\)-algebra \(A\) is called a \emph{PI algebra} if
there exists a nonzero polynomial
\(f\in\K\langle x_1,\ldots,x_n\rangle\) such that
\(f(a_1,\ldots,a_n)=0\) for all \(a_1,\ldots,a_n\in A\)
\cite[\S13.1]{mcconnell}. Such a polynomial is called a
\emph{polynomial identity} of \(A\).

We also recall that a connected graded algebra
\(A=\bigoplus_{n\geq0}A_n\), with \(A_0=\K\), is called
\emph{Artin--Schelter Gorenstein} of dimension \(d\) if it has finite
injective dimension \(d\) on both sides and there exists
\(\mathfrak l\in\mathbb Z\) such that
\[
\operatorname{Ext}_{A}^{i}({}_A\K,{}_AA)
\cong
\operatorname{Ext}_{A^{\mathrm{op}}}^{i}(\K_A,A_A)
\cong
\begin{cases}
0, & i\neq d,\\
\K(\mathfrak l), & i=d.
\end{cases}
\]
If, in addition, \(A\) has finite global dimension and finite
Gelfand--Kirillov dimension, then \(A\) is called
\emph{Artin--Schelter regular}
\cite[Definition~0.1]{ChanGaddisWonZhang2025}.

For an \(A\)-module \(M\), the \emph{homological grade} is defined by
\[
j_A(M)=\min\{i:\operatorname{Ext}_A^i(M,A)\neq0\},
\]
with \(j_A(M)=\infty\) if no such \(i\) exists. 

A Noetherian algebra \(A\) is said to satisfy the
\emph{Auslander condition} if, for every finitely generated left or
right \(A\)-module \(M\), every \(i\geq0\), and every submodule
\(N\) of \(\operatorname{Ext}_A^i(M,A)\), one has
\[
j_A(N)\geq i.
\]
The algebra \(A\) is called \emph{Auslander--regular} if it satisfies
the Auslander condition and has finite global dimension
\cite[Definition~2.1]{LezamaVenegas2017Homological}. Moreover, \(A\) is
called \emph{Cohen--Macaulay} with respect to the classical
Gelfand--Kirillov dimension if
\[
\operatorname{GKdim}(A)
=
j_A(M)+\operatorname{GKdim}(M)
\]
for every nonzero Noetherian \(A\)-module \(M\)
\cite[Definition~3.1]{LezamaVenegas2017Homological}.


 Moreover, a
\(\K\)-algebra \(A\) is called \emph{strongly Noetherian} if
\(A\otimes_{\K}C\) is Noetherian for every commutative Noetherian
\(\K\)-algebra \(C\).

Finally, let \(A\) be a graded algebra and let
\(A^e=A\otimes_{\K}A^{\operatorname{op}}\). The algebra \(A\) is called
\emph{skew Calabi--Yau} of dimension \(d\) if it is homologically smooth
and there exist an automorphism \(\mu\) of \(A\) and
\(\ell\in\mathbb Z\) such that
\[
\operatorname{Ext}_{A^e}^{i}(A,A^e)
\cong
\begin{cases}
0, & i\neq d,\\
{}^{1}\!A^{\mu}(\ell), & i=d,
\end{cases}
\]
as graded \(A^e\)-modules
\cite[Definition~0.1]{ReyesRogalskiZhang2014}.
A skew Calabi--Yau algebra is called \emph{Calabi--Yau} if its Nakayama
automorphism is inner. In the connected graded domain setting considered
here, this is equivalent to the Nakayama automorphism being the identity.

\section[The case where q is not a root of unity]{The ozone group }
The ozone group of \(U_q^{+}(B_2)\) is determined by distinguishing the
cases where \(q\) is and is not a root of unity. In the generic case, we
use the known descriptions of the automorphism group and the center. In
the root-of-unity case, with \(\ell=\operatorname{ord}(q^2)\), the
structure of the center allows us to restrict the automorphisms that fix
it pointwise and thereby determine the ozone group.


\subsection[The case where q is not a root of unity]{When \(q\) is not a root of unity}
Suppose that \(q\in\K^{*}\) is not a root of unity. In this case, the
automorphism group of \(U_q^{+}(B_2)\) is known. Andruskiewitsch and Dumas
\cite[Remark~3.2 and Proposition~3.3]{AndruskiewitschDumas2008}
observed that, in type \(B_2\), the Dynkin diagram automorphism group is
trivial and determined the subgroup of automorphisms stabilizing the ideal
generated by the central element \(z\). Subsequently, Launois
\cite[Proposition~4.3 and Theorem~4.4]{Launois2007} proved that every
automorphism of \(U_q^{+}(B_2)\) stabilizes this ideal. Consequently,
\[
\operatorname{Aut}_{\K}\!\left(U_q^{+}(B_2)\right)
=
\left\{
\psi_{\alpha,\beta}
\mid
\alpha,\beta\in\K^{*}
\right\}
\cong
(\K^{*})^{2},
\]
where
\[
\psi_{\alpha,\beta}(e_1)=\alpha e_1,\qquad
\psi_{\alpha,\beta}(e_2)=\beta e_2,\qquad
\psi_{\alpha,\beta}(e_3)=\alpha\beta\,e_3,\qquad
\psi_{\alpha,\beta}(z)=\alpha\beta^2\,z.
\]
Andruskiewitsch and Dumas
\cite[Lemma~3.1]{AndruskiewitschDumas2008} proved that
\(Z\!\left(U_q^{+}(B_2)\right)=\K[z,z']\), where
\[
z'=(1-q^{-4})(1-q^{-2})z_1,
\qquad
z_1=e_1e_2e_3+\frac{e_1z}{q^2-1}+\frac{e_3^2}{q^4-1}.
\]
Therefore,
\[
Z\!\left(U_q^{+}(B_2)\right)=\K[z,z_1].
\]
\begin{theorem}
Suppose that \(q\in\K^{*}\) is not a root of unity. Then
\[
\operatorname{Oz}\!\left(U_q^{+}(B_2)\right)
=
\{\operatorname{id},\tau\},
\]
where
\[
\tau(e_1)=e_1,
\qquad
\tau(e_2)=-e_2,
\qquad
\tau(e_3)=-e_3,
\qquad
\tau(z)=z.
\]
\end{theorem}

\begin{proof}
For
\(\psi_{\alpha,\beta}\in\operatorname{Aut}(U_q^{+}(B_2))\),
we have
\[
\psi_{\alpha,\beta}(z)=\alpha\beta^2z,
\qquad
\psi_{\alpha,\beta}(z_1)=\alpha^2\beta^2z_1.
\]
It follows that
\(\psi_{\alpha,\beta}\in\operatorname{Oz}\!\left(U_q^{+}(B_2)\right)\)
if and only if
\[
\alpha\beta^2=1,
\qquad
\alpha^2\beta^2=1.
\]
These equalities imply
\[
\alpha=1,
\qquad
\beta^2=1.
\]
Since \(\operatorname{char}(\K)=0\), we obtain \(\beta=\pm1\). Therefore,
\[
\operatorname{Oz}\!\left(U_q^{+}(B_2)\right)
=
\{\psi_{1,1},\psi_{1,-1}\}
=
\{\operatorname{id},\tau\}.
\]
\end{proof}
\subsection[The case where q is a root of unity]{When \(q\) is a root of unity}
Throughout this subsection, we assume that \(q\in\K^*\) is a primitive \(m\)-th root of unity, with \(m\geq5\). Let
$$
\ell=\operatorname{ord}(q^2)
=
\frac{m}{\gcd(m,2)}.
$$
In particular,
$$
\ell=
\begin{cases}
m, & \text{if \(m\) is odd},\\
m/2, & \text{if \(m\) is even}.
\end{cases}
$$
\\
For \(k\geq 0\), define the \(q\)-numbers
\[
[k]_{q^2}:=\frac{q^{2k}-1}{q^2-1},
\qquad
[k]_{q^{-4}}:=\frac{q^{-4k}-1}{q^{-4}-1}.
\]
which are well defined because \(m\geq5\) forces \(q^{2}\neq1\) and
\(q^{4}\neq1\). These \(q\)-numbers satisfy the recurrence relations
\[
q^{2k}+[k]_{q^2}=[k+1]_{q^2},
\qquad
q^{-4k}+[k]_{q^{-4}}=[k+1]_{q^{-4}},
\]
which will be used throughout the sequel.

\begin{lemma}{\cite[Lemma 2.1]{BeraMukherjee2026UqB2}} \label{lem:B2}
The identities below hold in  $U_q^+(B_2)$:
\begin{enumerate}[\rm (i)]
\item $e_2 e_3^{\,k}=q^{2k} e_3^{\,k} e_2+ [k]_{q^2}\, z\, e_3^{\,k-1}.$
  \item $e_2^{\,k} e_3=q^{2k} e_3 e_2^{\,k}+[k]_{q^2}\, z\, e_2^{\,k-1}.$
  \item $e_2 e_1^{\,k}=q^{-2k} e_1^{\,k} e_2-q^{-2}\,[k]_{q^{-4}}\, e_3 e_1^{\,k-1}.$
  \item $e_2^{k} e_1= q^{-2k} e_1 e_2^{k}- q^{-2} B_k\, e_3 e_2^{k-1}
- C_k\, z\, e_2^{k-2}$, 
where 
\begin{align*}
B_k=\frac{q^{2k}-q^{-2k}}{q^{2}-q^{-2}},
\;\;\;\;\;
C_k = q^{-2(k-1)}\,\frac{(q^{2k}-1)(q^{2(k-1)}-1)}{(q^{4}-1)(q^{2}-1)}.
\end{align*}
for $k \geq 2$.
\end{enumerate}
\end{lemma}
\begin{proposition}
{\cite[Corollary 2.2]{BeraMukherjee2026UqB2}}\label{centro2} 
The elements \(z\), \(e_{1}^{l}\), \(e_{2}^{l}\), and \(e_{3}^{l}\) are central in \(U_q^{+}(B_2)\).
\end{proposition}
\begin{proof}
It follows that
\begin{align*}
 q^{2l}=1,\quad [l]_{q^{2}}=0,\quad [l]_{q^{-4}}=0,\quad B_l=0,\quad C_l=0.
\end{align*}
From the defining relations of \(U_q^{+}(B_2)\), it follows that $z$ is a central element. On the other hand, by  Lemma \ref{lem:B2}--(i) with $k=l$ one obtains 
\[
e_2 e_3^{\,l}=q^{2l} e_3^{\,l} e_2 + [l]_{q^{2}} z e_3^{\,l-1}=e_3^{\,l} e_2.
\]
Moreover, from the relation $e_1 e_3=q^{-2} e_3 e_1$ it follows by induction that 
$e_1 e_3^{\,l}=q^{-2l} e_3^{\,l} e_1=e_3^{\,l} e_1$. Hence,  $e_3^{\,l}$ commute with $e_1$ and $e_2$.
Analogously, Lemma \ref{lem:B2}--(ii), for $k=l$, it follows
\[
e_2^{\,l} e_3 = q^{2l} e_3 e_2^{\,l} + [l]_{q^{2}} z e_2^{\,l-1}= e_3 e_2^{\,l}.
\]
Once again, Lemma \ref{lem:B2}--(iv) implies
\[
e_2^{\,l} e_1 = q^{-2l} e_1 e_2^{\,l} - q^{-2} B_l\, e_3 e_2^{\,l-1} - C_l\, z e_2^{\,l-2}
= e_1 e_2^{\,l};
\]
that is, $e_2^{\,l}$ commute with $e_1$ and $e_3$. From the relation $e_1 e_3=q^{-2} e_3 e_1$ can be proved that $e_3 e_1^{\,l}=q^{2l} e_1^{\,l} e_3 = e_1^{\,l} e_3$. Indeed, Lemma \ref{lem:B2}--(iii), assuming $k=l$, implies
\[
e_2 e_1^{\,l}= q^{-2l} e_1^{\,l} e_2 - q^{-2}[l]_{q^{-4}} e_3 e_1^{\,l-1}
= e_1^{\,l} e_2.
\]
Thus, $e_1^{\,l}$ commute  with $e_2$ and $e_3$. 
Finally, since  $e_1^{\,l}$, $e_2^{\,l}$, $e_3^{\,l}$ and $z$ commute with the generators of $U_q^{+}(B_2)$, it is concluded that these elements belong to its center.
\end{proof}

\begin{proposition}{\cite[Proposition 2.3]{BeraMukherjee2026UqB2}}\label{Propi}
 $U_q^{+}(B_2)$ is a PI algebra if and only if $q$ is a root of unity.
\end{proposition}
\begin{proof}
Consider the subalgebra $A:= \K\langle e_1,e_3\rangle$ of $U_q^{+}(B_2)$. The relation  $e_1e_3=q^{-2}e_3e_1$ shows that $A$ is isomorphic to the quantum plane with parameter $\lambda=q^{-2}$:
\[
A\ \cong\  \K\langle x,y\mid yx=\lambda xy\rangle,\qquad \lambda=q^{-2}.
\]
Hence, by \cite[Section \S 7.1]{DeConciniProcesi1993} one obtains  $q^{-2}$ is a root of unity, which implies that $q$ is a root of unity. Conversely, suppose that \(q\) is a root of unity and let
\(\ell=\operatorname{ord}(q^2)\). By
Proposition~\ref{centro2},

$$
R:=\K[e_1^\ell,e_2^\ell,e_3^\ell,z]
\subseteq Z(U_q^{+}(B_2)).
$$

For a monomial \(e_1^ae_2^be_3^c\), write
\(a=u\ell+i\), \(b=v\ell+j\), and \(c=w\ell+k\), with
\(0\leq i,j,k<\ell\). Using the identities in
Lemma~\ref{lem:B2} together with the relation
\(e_1e_3=q^{-2}e_3e_1\), the factors
\(e_1^\ell,e_2^\ell,e_3^\ell\) can be moved successively to the left, so that
$$
e_1^ae_2^be_3^c
\in
\sum_{0\leq i,j,k<\ell}
R\,e_1^ie_2^je_3^k.
$$
Hence, \(U_q^{+}(B_2)\) is finitely generated as an \(R\)-module, with
generating set
$$
\mathcal B=
\{e_1^ie_2^je_3^k:0\leq i,j,k<\ell\}.
$$
Since \(R\) is central, \cite[Corollary~13.1.13]{mcconnell} implies that
\(U_q^{+}(B_2)\) is PI.
\end{proof}
The central elements obtained above allow us to describe the center of \(U_q^{+}(B_2)\) completely. Indeed, Bera and Mukherjee \cite[Theorem~9.2]{BeraMukherjee2026UqB2} proved that it is generated by these elements together with an additional central element \(z_1\).
\begin{lemma}{\cite[Theorem~9.2]{BeraMukherjee2026UqB2}}
\label{centroU}
$$
Z\!\left(U_q^{+}(B_2)\right)
=
\K[e_1^\ell,e_2^\ell,e_3^\ell,z,z_1],
\qquad
z_1
=
e_1e_2e_3
+
\frac{e_1z}{q^2-1}
+
\frac{e_3^2}{q^4-1}.
$$
\end{lemma}
\begin{proposition}
\label{rigidez-ozono-UqB2}
Let \(\varphi\in\operatorname{Aut}(U_q^+(B_2))\) such that
\(\varphi(e_1^\ell)=e_1^\ell\) and \(\varphi(e_2^\ell)=e_2^\ell\). Then
there exist \(a,b\in\mu_\ell\) such that
$$
\varphi(e_1)=ae_1,\qquad
\varphi(e_2)=be_2,\qquad
\varphi(e_3)=ab\,e_3,\qquad
\varphi(z)=ab^2z.
$$
\end{proposition}
\begin{proof}
Let \(A=U_q^+(B_2)\). Consider the standard grading on \(A\), determined by

$$
\deg(e_1)=\deg(e_2)=1,\qquad
\deg(e_3)=2,\qquad
\deg(z)=3.
$$
The defining relations are homogeneous, so \(A\) is a connected
\(\mathbb N\)-graded algebra. Moreover, \(A\) admits a presentation as an
iterated Ore extension over \(\K[z,e_3]\), and hence it is a domain.
Consequently, \(\deg(x^\ell)=\ell\deg(x)\) for every \(x\neq0\). From
\(\varphi(e_1^\ell)=e_1^\ell\), we obtain
$$
\ell\deg(\varphi(e_1))
=\deg(\varphi(e_1)^\ell)
=\deg(e_1^\ell)
=\ell,
$$
and therefore \(\deg(\varphi(e_1))=1\). Since
\(A_0=\K\) and \(A_1=\K e_1\oplus\K e_2\), there exist \(a,b,c\in\K\) such that
$$
\varphi(e_1)=ae_1+be_2+c.
$$
Comparing the degree-zero components in
\(\varphi(e_1)^\ell=e_1^\ell\), we obtain \(c^\ell=0\), and hence \(c=0\). Thus,
$$
\varphi(e_1)=ae_1+be_2.
$$
Now consider the quotient \(A/\langle e_1\rangle\). Since \(e_1=0\), the relations

$$
e_2e_1=q^{-2}e_1e_2-q^{-2}e_3,
\qquad
e_2e_3=q^2e_3e_2+z
$$
imply that \(e_3=z=0\). Hence, \(A/\langle e_1\rangle\cong\K[e_2]\). Reducing the equality
$$
(ae_1+be_2)^\ell=e_1^\ell
$$
modulo \(\langle e_1\rangle\), we obtain \(b^\ell e_2^\ell=0\). Since
\(\K[e_2]\) is a domain, \(b=0\), and therefore \(\varphi(e_1)=ae_1\).
Moreover, from \(\varphi(e_1^\ell)=e_1^\ell\), we obtain \(a^\ell=1\), that is,
\(a\in\mu_\ell\).
\noindent Analogously, from \(\varphi(e_2^\ell)=e_2^\ell\) and the quotient
\(A/\langle e_2\rangle\cong\K[e_1]\), we obtain
$$
\varphi(e_2)=be_2
$$
for some \(b\in\mu_\ell\). Finally,
$$
\varphi(e_3)=ab(e_1e_2-q^2e_2e_1)=ab\,e_3
\qquad\text{and}\qquad
\varphi(z)=ab^2(e_2e_3-q^2e_3e_2)=ab^2z.
$$
\end{proof}
\begin{theorem}\label{ozono-UqB2}
$$
\operatorname{Oz}(U_q^+(B_2))\cong \mu_{\gcd(\ell,2)}.
$$
More precisely,
$$
\operatorname{Oz}(U_q^+(B_2))
\cong
\begin{cases}
\{1\},&\ell\text{ odd},\\[1mm]
\mu_2,&\ell\text{ even}.
\end{cases}
$$
\end{theorem}

\begin{proof}
Let \(\varphi\in\operatorname{Oz}(A)\). Since
\(e_1^\ell,e_2^\ell\in Z(A)\), by Proposition~\ref{rigidez-ozono-UqB2}
there exist \(a,b\in\mu_\ell\) such that
$$
\varphi(e_1)=ae_1,\qquad
\varphi(e_2)=be_2,\qquad
\varphi(e_3)=ab\,e_3,\qquad
\varphi(z)=ab^2z.
$$
Since \(z\in Z(U_q^{+}(B_2))\), we have \(\varphi(z)=z\), and hence
\(ab^2=1\). Likewise, since \(z_1\in Z(U_q^{+}(B_2))\),
$$
\varphi(z_1)
=a^2b^2\left(e_1e_2e_3+\frac{e_1z}{q^2-1}+\frac{e_3^2}{q^4-1}\right)
=a^2b^2z_1=z_1.
$$
Thus, \(a^2b^2=1\). Together with \(ab^2=1\), this implies \(a=1\) and
\(b^2=1\). Moreover, \(b^\ell=1\); therefore,
\(b\in\mu_{\gcd(\ell,2)}\). Consequently,
$$
\operatorname{Oz}(U_q^{+}(B_2))\subseteq
\{\psi_\varepsilon:\varepsilon\in\mu_{\gcd(\ell,2)}\}.
$$
Conversely, let \(\varepsilon\in\mu_{\gcd(\ell,2)}\). Then
\(\varepsilon^\ell=\varepsilon^2=1\). Consider the graded automorphism
$$
\psi_\varepsilon(e_1)=e_1,\qquad
\psi_\varepsilon(e_2)=\varepsilon e_2.
$$
Then \(\psi_\varepsilon(e_3)=\varepsilon e_3\) and
\(\psi_\varepsilon(z)=z\). Moreover,
$$
\psi_\varepsilon(e_1^\ell)=e_1^\ell,\qquad
\psi_\varepsilon(e_2^\ell)=e_2^\ell,\qquad
\psi_\varepsilon(e_3^\ell)=e_3^\ell,\qquad
\psi_\varepsilon(z_1)=\varepsilon^2z_1=z_1.
$$
Since
\(Z(U_q^{+}(B_2))=\K[e_1^\ell,e_2^\ell,e_3^\ell,z,z_1]\),
\(\psi_\varepsilon\) fixes \(Z(U_q^{+}(B_2))\) pointwise. Therefore,
$$
\psi_\varepsilon\in\operatorname{Oz}(U_q^{+}(B_2)).
$$
\end{proof}


\begin{proposition}\label{prop:normal-explicito}
Suppose that \(\ell\) is even. Then \(e_1^{\ell/2}\) is a noncentral normal
element of \(U_q^{+}(B_2)\), and
\[
e_1^{\ell/2}\,e_1=e_1\,e_1^{\ell/2},\qquad
e_1^{\ell/2}\,e_2=-\,e_2\,e_1^{\ell/2},\qquad
e_1^{\ell/2}\,e_3=-\,e_3\,e_1^{\ell/2},\qquad
e_1^{\ell/2}\,z=z\,e_1^{\ell/2}.
\]
Consequently \(\eta_{e_1^{\ell/2}}=\psi_{-1}\) generates
\(\operatorname{Oz}(U_q^{+}(B_2))\).
\end{proposition}

\begin{proof}
Put \(k=\ell/2\). From \(e_1e_3=q^{-2}e_3e_1\) one gets by induction
\(e_1^{k}e_3=q^{-2k}e_3e_1^{k}\); since \(q^{2}\) has order \(\ell\) and
\(0<k<\ell\), the scalar \(q^{-2k}=(q^{2})^{-\ell/2}\) is a square root of
\(1\) different from \(1\), that is, \(q^{-2k}=-1\). Hence
\(e_1^{k}e_3=-e_3e_1^{k}\).

By Lemma~\ref{lem:B2}(iii),
\[
e_2 e_1^{k}=q^{-2k} e_1^{k} e_2-q^{-2}\,[k]_{q^{-4}}\, e_3 e_1^{k-1},
\qquad
[k]_{q^{-4}}=\frac{q^{-4k}-1}{q^{-4}-1}=\frac{q^{-2\ell}-1}{q^{-4}-1}=0,
\]
because \(q^{2\ell}=1\) and \(q^{4}\neq1\). Therefore
\(e_2e_1^{k}=-e_1^{k}e_2\). Since \(z\) is central, \(e_1^{k}\) commutes with
\(z\) and with \(e_1\), so \(e_1^{k}U_q^{+}(B_2)=U_q^{+}(B_2)e_1^{k}\) and
\(e_1^{k}\) is normal; it is not central because \(\ell\geq3\) forces
\(-1\neq1\). Finally, with the convention
\(\eta_a(x)=a^{-1}xa\) used in \cite[Lemma~1.10]{ChanGaddisWonZhang2025}, the
relations above give \(\eta_{e_1^{k}}(e_1)=e_1\),
\(\eta_{e_1^{k}}(e_2)=-e_2\), \(\eta_{e_1^{k}}(e_3)=-e_3\) and
\(\eta_{e_1^{k}}(z)=z\), that is, \(\eta_{e_1^{k}}=\psi_{-1}\).
\end{proof}




\begin{corollary}
\label{cor:rango-ozono-UqB2}
$$
\operatorname{rk}_{Z(U_q^{+}(B_2))}(U_q^{+}(B_2))=\ell^2
\qquad\text{and}\qquad
|\operatorname{Oz}(U_q^{+}(B_2))|
<
\operatorname{rk}_{Z(U_q^{+}(B_2))}(U_q^{+}(B_2)).
$$
\end{corollary}





\begin{proof}
Since \(U_q^{+}(B_2)\) is prime, PI, and finitely generated as a module over its center, and by \cite[Theorem~2.7]{BeraMukherjee2026UqB2}
$$
\operatorname{PIdeg}(U_q^{+}(B_2))=\ell,
$$

 Let
\(Q\) denote its total ring of fractions. Since \(U_q^{+}(B_2)\) is PI, one has
\(\GKdim(U_q^{+}(B_2))=\GKdim(Z(U_q^{+}(B_2)))\), so the hypothesis
\(\GKdim(A)<\GKdim(Z(A))+1\) of \cite[Theorem~2.1]{LezamaVenegas2020Serdica} is
satisfied and therefore \(Z(Q)\cong Q(Z(U_q^{+}(B_2)))\). By Posner's theorem
\cite[Theorem~13.6.5]{mcconnell}, \(Q\) is a central simple algebra over
\(Z(Q)\) of degree \(\operatorname{PIdeg}(U_q^{+}(B_2))=\ell\), whence

$$
\operatorname{rk}_{Z(U_q^{+}(B_2))}U_q^{+}(B_2) = \operatorname{PIdeg}(U_q^{+}(B_2))^2 = \ell^2.
$$
Moreover, by Theorem~\ref{ozono-UqB2},
$$
|\operatorname{Oz}(U_q^{+}(B_2))|
=
\gcd(\ell,2).
$$
Since \(m\geq5\), we have \(\ell\geq3\), and therefore
$$
\gcd(\ell,2)
=
|\operatorname{Oz}(U_q^{+}(B_2))|
<
\ell^2
=
\operatorname{rk}_{Z(U_q^{+}(B_2))}(U_q^{+}(B_2)).
$$
\end{proof}











\section{Some homological properties of
\texorpdfstring{$U_q^{+}(B_2)$}{Uq+(B2)}}
Graded skew PBW extensions were introduced by Suárez
\cite{Suarez2017KoszulityGradedSkewPBW} in 2017 as a generalization of
graded iterated Ore extensions. Subsequently, the first author and Suárez
\cite{GomezSuarez2020} studied their relationship with double Ore
extensions and proved that every graded skew PBW extension in two
variables over an Artin--Schelter regular algebra is again
Artin--Schelter regular. In what follows, we show that
\(U_q^{+}(B_2)\) can be realized as a graded skew PBW extension in two
variables and, as a consequence, establish that \(U_q^{+}(B_2)\) is
Artin--Schelter regular. We also investigate whether this algebra is Auslander-regular, 
Cohen--Macaulay, strongly noetherian and skew Calabi--Yau.
\begin{definition}[{\cite[Definition~1]{GallegoLezama}}]
Let \(R\subseteq A\) be rings. We say that \(A\) is a
\emph{skew PBW extension} of \(R\), and write
$$
A=\sigma(R)\langle x_1,\ldots,x_n\rangle,
$$
if \(A\) is a free left \(R\)-module with basis
$$
\operatorname{Mon}(A)=
\{x_1^{\alpha_1}\cdots x_n^{\alpha_n}:
(\alpha_1,\ldots,\alpha_n)\in\mathbb N^n\},
$$
and there exist injective endomorphisms \(\sigma_i:R\to R\),
\(\sigma_i\)-derivations \(\delta_i:R\to R\), and elements
\(c_{i,j}\in R\setminus\{0\}\) such that
$$
x_ir=\sigma_i(r)x_i+\delta_i(r),\qquad
x_jx_i-c_{i,j}x_ix_j\in R+Rx_1+\cdots+Rx_n,
$$
for \(r\in R\) and \(1\leq i<j\leq n\).
The extension is called \emph{bijective} if each \(\sigma_i\) is an
automorphism and each \(c_{i,j}\) is invertible in \(R\).
\end{definition}

\begin{definition}[{\cite[Proposition~2.7]{Suarez2017KoszulityGradedSkewPBW}}]
Let \(R=\bigoplus_{m\geq0}R_m\) be an \(\mathbb N\)-graded algebra and
let \(A=\sigma(R)\langle x_1,\ldots,x_n\rangle\) be a bijective skew PBW
extension of \(R\). We say that \(A\) is a \emph{graded skew PBW
extension} if \(\deg(x_i)=1\), each \(\sigma_i\) preserves the grading,
\(\delta_i(R_m)\subseteq R_{m+1}\), and, for \(1\leq i<j\leq n\),
$$
x_jx_i-c_{i,j}x_ix_j\in R_2+R_1x_1+\cdots+R_1x_n,
\qquad c_{i,j}\in R_0.
$$
In this case,
$$
A_p=\operatorname{span}_{\K}
\{r_tx_1^{\alpha_1}\cdots x_n^{\alpha_n}:
r_t\in R_t,\ t+\alpha_1+\cdots+\alpha_n=p\}.
$$
\end{definition}
\begin{proposition}
\label{prop:PBW-graduada-UqB2}
Let \(q\in\K^*\). Then \(U_q^{+}(B_2)\) is a graded bijective skew PBW extension
in two variables over \(R=\K[z,e_3]\), where \(\deg(e_3)=2\),
\(\deg(z)=3\), and \(\deg(e_1)=\deg(e_2)=1\).
\end{proposition}
\begin{proof}
Let \(R=\K[z,e_3]\), endowed with the grading determined by
\(\deg(e_3)=2\) and \(\deg(z)=3\). By
\cite[Section~3.1.2]{AndruskiewitschDumas2008}, the monomials
$$
\{z^ie_3^je_1^ke_2^r:i,j,k,r\geq0\}
$$
form a PBW basis of \(U_q^{+}(B_2)\). Hence,
\(U_q^{+}(B_2)\) is a free left \(R\)-module with basis
\(\{e_1^ke_2^r:k,r\geq0\}\).
The relations of \(e_1,e_2\) with \(R\) are determined by
\(\sigma_1(z)=z\), \(\sigma_1(e_3)=q^{-2}e_3\),
\(\sigma_2(z)=z\), \(\sigma_2(e_3)=q^2e_3\),
\(\delta_1=0\), \(\delta_2(z)=0\), and \(\delta_2(e_3)=z\), together with
$$
e_2e_1-q^{-2}e_1e_2=-q^{-2}e_3\in R_2.
$$
Thus, \(U_q^{+}(B_2)\) is a skew PBW extension of \(R\).
Moreover, \(\sigma_1,\sigma_2\) are automorphisms and
\(q^{-2}\in R_0^*\), so the extension is bijective. Finally,
\(\deg(e_1)=\deg(e_2)=1\), the maps \(\sigma_1,\sigma_2\) preserve the
grading, \(\delta_1=0\) is compatible with the grading, and \(\delta_2\)
has degree \(1\), since
\(\deg(\delta_2(e_3))=3=\deg(e_3)+1\) and \(\delta_2(z)=0\).
Furthermore, the relation above satisfies the graded condition because
its remaining term belongs to \(R_2\) and \(q^{-2}\in R_0\). Therefore,
\(U_q^{+}(B_2)=\sigma(R)\langle e_1,e_2\rangle\) is a graded skew PBW
extension of \(R\).
\end{proof}
\begin{theorem}
\label{AS-regular-UqB2}
Let \(q\in\K^*\). Then \(U_q^{+}(B_2)\) is an
Artin--Schelter regular algebra of global dimension \(4\).
\end{theorem}
\begin{proof}
By Proposition~\ref{prop:PBW-graduada-UqB2},
\(U_q^{+}(B_2)\) is a graded skew PBW extension in two variables
over \(R=\K[z,e_3]\), where \(\deg(e_3)=2\) and \(\deg(z)=3\).
Since \(R\) is a commutative graded polynomial ring in two homogeneous
variables of positive degrees, it is Artin--Schelter regular of global
dimension \(2\). By \cite[Theorem~3.10]{GomezSuarez2020},
\(U_q^{+}(B_2)\) is Artin--Schelter regular and
$$\operatorname{gldim}(U_q^{+}(B_2))=\operatorname{gldim}(R)+2=4.$$
\end{proof}

\begin{remark}
As noted in the Introduction, \(U_q^{+}(B_2)\) admits a presentation as
a graded iterated Ore extension. Since the corresponding automorphisms
preserve the grading and the derivation is homogeneous of degree \(1\),
successive applications of
\cite[Proposition~2]{ArtinSchelterTate1991Quantum} also show that
\(U_q^{+}(B_2)\) is Artin--Schelter regular of global dimension \(4\).

\noindent Thus, the iterated Ore extension presentation provides an alternative
and direct proof of the previous theorem. However, realizing
\(U_q^{+}(B_2)\) as a bijective skew PBW extension allows one to apply
systematically a variety of transfer results from the base algebra to the
extension. In particular, the theory of skew PBW extensions provides
general criteria for preserving properties such as Auslander regularity,
the Cohen--Macaulay property, and strong Noetherianity
\cite{LezamaVenegas2017Homological}, as well as other homological,
structural, and \(K\)-theoretic results for bijective skew PBW extensions
\cite{LezamaReyes2014Homological}.
\end{remark}




\begin{theorem}
\label{teo:UqB2-Auslander-CM}
Let \(q\in\K^*\). Then \(U_q^{+}(B_2)\) is Auslander--regular, Cohen--Macaulay, and strongly Noetherian.
\end{theorem}


\begin{proof}
The algebra \(R=\K[z,e_3]\) is a commutative polynomial ring; in
particular, it is Auslander--regular, Cohen--Macaulay, and strongly
Noetherian. By Proposition~\ref{prop:PBW-graduada-UqB2},
\(U_q^{+}(B_2)=\sigma(R)\langle e_1,e_2\rangle\) is a bijective skew PBW
extension of \(R\). Hence, by
\cite[Theorem~2.9]{LezamaVenegas2017Homological},
\(U_q^{+}(B_2)\) is Auslander--regular.

Moreover, \(R\) is connected graded and
\(\sigma_i(R_n)\subseteq R_n\) for every \(n\geq0\) and \(i=1,2\).
Therefore, \cite[Theorem~3.9]{LezamaVenegas2017Homological} implies that
\(U_q^{+}(B_2)\) is Cohen--Macaulay. Finally, since \(R\) is strongly
Noetherian, \cite[Theorem~4.3]{LezamaVenegas2017Homological} implies that
\(U_q^{+}(B_2)\) is strongly Noetherian.
\end{proof}

Reyes, Rogalski, and
Zhang also proved that a connected graded algebra is skew
Calabi--Yau if and only if it is Artin--Schelter regular
\cite[Lemma~1.2]{ReyesRogalskiZhang2014}.
\begin{theorem}
\label{teo:UqB2-skew-CY}
Let \(q\in\K^*\). Then \(U_q^{+}(B_2)\) is skew Calabi--Yau
of dimension \(4\).
\end{theorem}
\begin{proof}
By Theorem~\ref{AS-regular-UqB2},
\(U_q^{+}(B_2)\) is Artin--Schelter regular of global dimension \(4\).
Moreover, with the grading considered above,
\(U_q^{+}(B_2)\) is a connected graded algebra. Therefore, by
\cite[Lemma~1.2]{ReyesRogalskiZhang2014},
\(U_q^{+}(B_2)\) is skew Calabi--Yau of dimension \(4\).
\end{proof}

\subsection{Some applications of the ozone group computation}
In the previous sections, we established that, when \(q\) is a root of
unity, \(U_q^{+}(B_2)\) is a PI algebra, Artin--Schelter regular, Auslander regular, and
strongly Noetherian, and hence Noetherian. We also determined its ozone
group explicitly:
$$
\operatorname{Oz}(U_q^{+}(B_2))
\cong \mu_{\gcd(\ell,2)},
\qquad \ell=\operatorname{ord}(q^2).
$$
These properties place \(U_q^{+}(B_2)\) within the framework considered
by Chan, Gaddis, Won, and Zhang \cite{ChanGaddisWonZhang2025}. By combining
their general results with the ozone group computation obtained above,
we derive below several consequences concerning the Calabi--Yau property,
the Gorenstein property of the center, and the normal elements of
\(U_q^{+}(B_2)\).  Here, Gorenstein is understood in the usual commutative sense for \(Z(U_q^{+}(B_2))\).

\begin{theorem}
\label{teo:UqB2-center-Gorenstein}
Let \(q\in\K^*\) be a primitive \(m\)-th root of unity, with
\(m\geq5\), and let

$$
\ell=\operatorname{ord}(q^2)=\frac{m}{\gcd(m,2)}.
$$

If \(\ell\) is odd, then \(U_q^{+}(B_2)\) is Calabi--Yau and
\(Z(U_q^{+}(B_2))\) is Gorenstein.
\end{theorem}

\begin{proof}
By Proposition~\ref{Propi}, \(U_q^{+}(B_2)\) is a PI algebra; by
Theorem~\ref{AS-regular-UqB2}, it is Artin--Schelter regular; and by
Theorem~\ref{teo:UqB2-Auslander-CM}, it is strongly Noetherian, hence
Noetherian. Moreover, since \(\ell\) is odd,
Theorem~\ref{ozono-UqB2} implies that
$$
\operatorname{Oz}(U_q^{+}(B_2))=\{1\}.
$$
Therefore, by \cite[Theorem~1.11]{ChanGaddisWonZhang2025},
\(U_q^{+}(B_2)\) is Calabi--Yau and, consequently,
\(Z(U_q^{+}(B_2))\) is Gorenstein.
\end{proof}
\begin{proposition}
\label{prop:normales-ozono-UqB2}
Let \(N^*\) be the multiplicative monoid of nonzero normal elements of
\(U_q^{+}(B_2)\), and let \(\sim\) be the equivalence relation defined by
$$
a\sim b
\quad\Longleftrightarrow\quad
z_1a=z_2b
$$
for some
\(z_1,z_2\in Z(U_q^{+}(B_2))\setminus\{0\}\). Then
$$
N^*/{\sim}
\cong
\operatorname{Oz}(U_q^{+}(B_2))
\cong
\mu_{\gcd(\ell,2)}.
$$
In particular, if \(\ell\) is odd, every normal element of
\(U_q^{+}(B_2)\) is central; if \(\ell\) is even, the class of \(e_1^{\ell/2}\) generates
\(N^*/{\sim}\).
\end{proposition}
\begin{proof}
In the proof of \cite[Lemma~1.10]{ChanGaddisWonZhang2025}, the authors
consider the surjective map
$$
N^*\longrightarrow\operatorname{Oz}(U_q^{+}(B_2)),
\qquad
a\longmapsto\eta_a,\quad \eta_a(x)=a^{-1}xa,
$$
whose surjectivity follows from
\cite[Lemma~1.9]{ChanGaddisWonZhang2025}. Moreover, for the equivalence
relation \(a\sim b\) whenever \(z_1a=z_2b\) for some
\(z_1,z_2\in Z(U_q^{+}(B_2))\setminus\{0\}\), one has
$$
N^*/{\sim}\cong\operatorname{Oz}(U_q^{+}(B_2)).
$$
By Theorem~\ref{ozono-UqB2},
$$
N^*/{\sim}\cong\operatorname{Oz}(U_q^{+}(B_2))
\cong\mu_{\gcd(\ell,2)}.
$$
If \(\ell\) is odd, \cite[Lemma~1.10]{ChanGaddisWonZhang2025}
implies that every normal element of \(U_q^{+}(B_2)\) is central; If \(\ell\) is even, Proposition
\ref{prop:normal-explicito} shows that \(e_1^{\ell/2}\) is normal and
noncentral with \(\eta_{e_1^{\ell/2}}=\psi_{-1}\), which generates
\(\operatorname{Oz}(U_q^{+}(B_2))\); hence its class generates
\(N^*/{\sim}\).
\end{proof}

\begin{question*}
A closely related generalization of \(U_q^{+}(B_2)\) is the algebra
\(U_{r,s}^{+}(B_2)\). When \(r\) and \(s\) are roots of unity, its PI
property, PI degree, and simple modules have recently been studied
\cite{MukherjeePandey2025}, and distinguished normal elements also arise
in this analysis. These features suggest studying its ozone group.
\end{question*}

\subsection*{Future work}
A natural direction for future work is to extend the study of ozone groups to other families of quantum algebras at roots of unity. In recent years,
Sanu Bera, Snehashis Mukherjee, and Sugata Mandal, together with other
collaborators, have investigated several such families, studying
properties such as the PI condition, PI degree, the existence of central
and normal elements, and the structure of their simple modules. These
families include the two-parameter quantum Heisenberg algebras
\(\mathcal H_{p,q}\), quantized Weyl algebras, the bialgebra \(M(p,q)\),
and the Dipper--Donkin and reflection equation quantum matrix algebras
\cite{BeraMandalNandy2024, BeraMukherjee2024DipperDonkin,Mukherjee2026Bialgebra,  MukherjeeBera2025}. These results provide part of the structural information that naturally
arises in the study of ozone groups, and it would therefore be interesting
to determine this invariant for those families for which it is not yet
known.

\section*{Acknowledgments}

We thank the members of the Research Seminar on Noncommutative Algebra for
the discussions, comments, and suggestions that contributed to the development
of this work. Special thanks are due to Professor Héctor Julio Suárez Suárez
for his active participation, guidance, and continued support throughout this
research.

%%%%%%%%%%%%%%%%%%%%%%%%%%%%%%%%%%%%%%%%%%%%%%%%%
\bibliographystyle{plain}
\bibliography{biblio}
%%%%%%%%%%%%%%%%%%%%%%%%%%%%%%%%%%%%%%%%%%%
\bigskip
\noindent
James Yair Gómez Lozano.:\\
Grupo de Álgebra y Análisis,\\
Universidad Pedagógica y Tecnológica de Colombia,\\
Tunja, Colombia\\
james.gomez@uptc.edu.co\\

Helbert Javier Venegas Ramírez.:\\
Universidad Militar Nueva Granada, Bogot\'a, Colombia,\\
Bogotá, Colombia\\
Helbert.venegas@unimilitar.edu.co



\end{document}



